% ==============================================================================
% Chapter 11: Interlude---Exotic Modular Morphisms and Isomorphisms
% ==============================================================================
\chapter{Interlude---Exotic Modular Morphisms and Isomorphisms}
\label{ch:11}

\section{Motivation}
\label{sec:11.1}

\kmitem\label{item:11.1.1}
Let $R$ be a ring, $\mathcal{P}$ and $\mathcal{P}'$ two moduli problems on $(\mathrm{Ell}/R)$. So far we have been content to always view $\mathcal{P}$ and $\mathcal{P}'$ as lying over $(\mathrm{Ell}/R)$, and we have only considered morphisms between them which are morphisms of moduli problems on $(\mathrm{Ell}/R)$ (i.e., morphisms between functors on $(\mathrm{Ell}/R)$). However, it is convenient, and often essential, to view $\mathcal{P}$ and $\mathcal{P}'$ as ``abstract'' objects, rather than as ``objects over $(\mathrm{Ell}/R)$'', and to have a notion of ``abstract'' morphism between them. For example, if $\mathcal{P}$ and $\mathcal{P}'$ are both representable, by universal families $E/\mathfrak{M}(\mathcal{P})$ and $E'/\mathfrak{M}(\mathcal{P}')$, then we certainly want to be able to speak of $R$-maps between the $R$-schemes $\mathfrak{M}(\mathcal{P})$ and $\mathfrak{M}(\mathcal{P}')$.

\section[``Abstract'' morphisms]{``Abstract'' morphisms\protect\footnotemark}
\footnotetext{Interchangeably called ``exotic.''}
\label{sec:11.2}

\kmitem\label{item:11.2.1}
With this preamble, we now define an ``abstract'' morphism
\[
\mathcal{P} \to \mathcal{P}'
\]
to be a rule which for every $R$-scheme $S$ and every elliptic $E/S$ with level $\mathcal{P}$-structure $x \in \mathcal{P}(E/S)$, assigns an elliptic curve $E'/S$ with level $\mathcal{P}'$-structure $x' \in \mathcal{P}'(E'/S)$, and which to every morphism in $(\mathrm{Ell}/R)$
\[
\begin{tikzcd}[row sep=1.8em, column sep=1.8em]
x_1 = (\alpha, f)^*(x) \in \mathcal{P}(E_1/S_1);\ E_1 \arrow[r, "\alpha"] \arrow[d] & E;\ x \in \mathcal{P}(E/S) \arrow[d] \\
S_1 \arrow[r, "f"] & S
\end{tikzcd}
\]
assigns a morphism in $(\mathrm{Ell}/R)$ with the same $f$
\[
\begin{tikzcd}[row sep=1.8em, column sep=1.8em]
(x_1)' \in \mathcal{P}'(E_1'/S_1);\ E_1' \arrow[r, "\alpha'"] \arrow[d] & E';\ x' \in \mathcal{P}'(E'/S) \arrow[d] \\
S_1 \arrow[r, "f"] & S
\end{tikzcd}
\]
such that $(x_1)' = (\alpha', f)^*(x')$, and such that the construction
\[
(\alpha, f) \mapsto (\alpha', f)
\]
is compatible with composition of morphisms in $(\mathrm{Ell}/R)$.

Clearly if both $\mathcal{P}$ and $\mathcal{P}'$ are representable, then any such abstract morphism induces a morphism of $R$-schemes
\[
\mathfrak{M}(\mathcal{P}) \to \mathfrak{M}(\mathcal{P}') ,
\]
simply by passing to isomorphism classes.

\section{Some basic examples}
\label{sec:11.3}

\kmitem\label{item:11.3.1}
The involution ``$W_N$'' of $[\Gamma_0(N)]$ is defined by
\[
\begin{tikzcd}[row sep=1.8em]
(E/S, \text{cyclic subgroup } G \text{ of order } N) \arrow[d, mapsto] \\
(E' = (E \bmod G)/S,\ G' = E[N] \bmod G) .
\end{tikzcd}
\]

\kmitem\label{item:11.3.2}
The automorphism ``$W_N$'' of $([\mathrm{bal.}\ \Gamma_1(N)]^{\mathrm{can}}, [\Gamma_1(M)])$ with $(N,M) = 1$. Let us consider an $E/S$ with both a balanced $\Gamma_1(N)$-structure
\[
P;\ E \kmdual{\pi}{\pi^t} E';\ Q
\]
and a $[\Gamma_1(M)]$-structure $A \in E[M](S)$. Because $(N,M) = 1$, $\pi$ and $\pi^t$ each induces isomorphisms $E[M] \kmdual{\pi}{\pi^t} E'[M]$, so that $\pi A$ has exact order $M$ on $E'/S$. Then ``$W_N$'' is the automorphism
\[
\begin{tikzcd}[row sep=1.8em]
P;\ E \kmdual{\pi}{\pi^t} E';\ Q;\ A \arrow[d, mapsto] \\
-Q;\ E' \kmdual{\pi}{\pi^t} E;\ P;\ \pi A .
\end{tikzcd}
\]
The square $W_N^2$ of this exotic automorphism is the automorphism of $([\mathrm{bal.}\ \Gamma_1(N)]^{\mathrm{can}}, [\Gamma_1(M)])$ as moduli problem on $(\mathrm{Ell}/\Z[\zeta_N])$ which is
\[
(P, Q, A) \mapsto (-P, -Q, NA) .
\]

\kmitem\label{item:11.3.3}
The two projections of $[\mathrm{bal.}\ \Gamma_1(p^n)]^{\mathrm{can}}$ to\linebreak $[\mathrm{bal.}\ \Gamma_1(p^{n-1})]^{\mathrm{can}}$.

The first is:
\[
\begin{tikzcd}[row sep=1.8em]
P_0;\ E_0 \kmdual{\pi_{0,1}}{\lambda_{1,0}} E_1 \rightleftarrows \cdots \kmdual{\pi_{n-1,n}}{\lambda_{n,n-1}} E_n;\ Q_n \arrow[d, mapsto] \\
pP_0;\ E_0 \rightleftarrows \cdots \rightleftarrows E_{n-1};\ \lambda_{n,n-1}(Q_n) .
\end{tikzcd}
\]
The second is:
\[
\begin{tikzcd}[row sep=1.8em]
\text{--- --- --- --- --- --- --- --- ---} \arrow[d, mapsto] \\
\pi_{0,1}(P_0);\ E_1 \rightleftarrows \cdots \rightleftarrows E_n;\ pQ_n .
\end{tikzcd}
\]

\kmitem\label{item:11.3.4}
The isomorphism
\[
[\Gamma_0(p^{2n}); n, n]^{\mathrm{can}} \xrightarrow{\ \sim\ } [\Gamma(p^n)]^{\mathrm{can}}
\]
defined by
\[
\begin{tikzcd}[row sep=1.2em, column sep=small]
& P_n \\
& \text{\rotatebox{-90}{$\in$}} \\
E_0 \kmdual{\pi_{0,n}}{\lambda_{n,0}} & E_n \kmdual{\pi_{n,2n}}{\lambda_{2n,n}} & E_{2n} \\
& \text{\rotatebox{90}{$\in$}} \\
& Q_n \arrow[d, mapsto] \\
& E_n;\ (P_n, Q_n) .
\end{tikzcd}
\]
To see that the map is defined, we must show that $(P_n, Q_n)$ is a Drinfeld $p^n$-basis on $E_n$. This is clear if $p$ is invertible (for example, because $e_{p^n}(P_n, Q_n) = \langle P_n, \pi_{n,2n} Q_n \rangle$ is a primitive $p^n$'th root of unity), and therefore holds universally, because $[\Gamma_0(p^{2n}); n, n]$ is flat over $\Z$.

The inverse map is defined by
\[
\resizebox{\linewidth}{!}{$\displaystyle
\begin{array}{c}
E_n;\ (P_n, Q_n) \\[4pt]
\Big\downarrow \\[6pt]
\begin{array}{c} E_n/\text{cyclic } p^n\text{-group} = E_0 \\[-3pt] \text{generated by } Q_n \end{array}
\xleftarrow{\ \bmod Q_n\ } E_n \xrightarrow{\ \bmod P_n\ }
E_{2n} = \begin{array}{c} E_n/\text{cyclic } p^n\text{-group} \\[-3pt] \text{generated by } P_n . \end{array}
\end{array}
$}
\]
To see that this map is well defined, we must see that the lower diagram is cyclic in standard order. Since the question is f.p.p.f.\ local on $S$, we may suppose given $P_0$ on $E_0$ such that $\pi_{0,n} P_0 = P_n$. We must show that $p^n P_0$ generates $\Ker\pi_{0,n}$. We easily calculate
\[
p^n P_0 = \lambda_{n,0}\pi_{0,n} P_0 = \lambda_{n,0} P_n = P_n \bmod Q_n ,
\]
and so we are reduced to the already noted fact (\ref{thm:5.5.2}) that
\[
(P_n \bmod Q_n;\ E_0 \rightleftarrows E_n;\ Q_n)
\]
is indeed a balanced $\Gamma_1(p^n)$-structure.

\kmitem\label{item:11.3.5}
The isomorphism
\[
[\Gamma(p^n)]/\begin{pmatrix} * & 0 \\ 0 & * \end{pmatrix} \xrightarrow{\ \sim\ } [\Gamma_0(p^{2n})] .
\]
We have already noted that $[\Gamma(p^n)]$ modulo the full Cartan subgroup of $\GL(2, \Z/p^n\Z)$ is regular and two-dimensional. It is easy to check that its modular interpretation is this: a $[\Gamma(p^n)]/\begin{pmatrix} * & 0 \\ 0 & * \end{pmatrix}$-structure on $E/S$ is a pair of cyclic subgroups of order $p^n$ in $E/S$, say $G_1$ and $G_2$, which are \emph{transverse} in the sense that locally f.p.p.f., if $P$ is a generator of $G_1$ and if $Q$ is a generator of $G_2$, then $(P,Q)$ is a Drinfeld $p^n$-basis of $E$.

This said, the asserted isomorphism is given by
\[
\begin{tikzcd}[row sep=1.8em]
E;\ G_1, G_2 \arrow[d, mapsto] \\
E \bmod G_1 \longleftarrow E \longrightarrow E \bmod G_2 ,
\end{tikzcd}
\]
and its inverse is given by
\[
\begin{tikzcd}[row sep=1.8em]
E_0 \kmdual{\pi_{0,n}}{\lambda_{n,0}} E_n \kmdual{\pi_{n,2n}}{\lambda_{2n,n}} E_{2n} \arrow[d, mapsto] \\
E_n;\ \Ker(\lambda_{n,0}), \Ker(\pi_{n,2n}) .
\end{tikzcd}
\]